\section{Method}\label{sec:method}
\pending{This section describes the design as specified before any experiment; parameters are the defaults of the
pre-registered starting configuration and may change in the capped tuning phase (Section~\ref{sec:design}).}

\subsection{Problem form, violation and feasibility}
We consider $\min\, c^\top x$ subject to $r^l \le Ax \le r^u$, $l \le x \le u$ and $x_j \in \mathbb{Z}$ for
$j \in I$, read from an MPS file; maximisation is handled through a sign. The population compares \emph{keys}
(lower is better). The total violation of an individual is
\[
V(x) = \sum_i \frac{\max(r^l_i - a_i, 0) + \max(a_i - r^u_i, 0)}{s_i} + \sum_j (\text{bound violation})_j
     + \sum_{j \in I} |x_j - \operatorname{round}(x_j)|, \qquad a = Ax,
\]
with a row scale $s_i$ (the row's 2-norm by default). $V$ orders infeasible individuals; it is never the feasibility
test. An individual is \emph{feasible} if every row, bound and integrality condition holds within a tolerance under the
partner solver's own rule --- relative to $\max(|a_i|, |r_i|, 1)$ for SCIP, absolute for HiGHS --- tested row by row,
never as a sum. Objective values are always computed in double precision. Every solution passed to a solver is
re-checked on the CPU in double precision first, and every solution that enters a reported metric is re-checked by an
independent exact checker.

\subsection{The three ranks}
Each generation, the candidate set (population, offspring and solutions injected from the solver) is reduced to the
population size $P$:
\begin{enumerate}
\item Exact duplicates are removed (their provenance is merged into the copy that is kept).
\item The incumbent key $k^*$ is the better of the best feasible key in the candidate set and the solver's incumbent.
\item \textbf{Rank 1 (feasible good):} the best $\lceil y \cdot \min(\#\text{feasible}, P) \rceil$ feasible candidates
  ($y = 10\,\%$) are protected; candidates with equal keys leave Rank~1.
\item \textbf{Rank 0 (infeasible good):} infeasible candidates with key below $k^*$ (by a relative margin); those beyond
  the solver's proven bound are removed. At most $\lfloor x P \rfloor$ survive ($x = 10\,\%$): over the cap, the
  non-dominated fronts in (key, $V$) are taken in order and the overflowing front is thinned by repeatedly removing the
  point nearest to another. Candidates over the cap are removed, not demoted, so $x$ bounds the number of super-optimal
  infeasible survivors.
\item \textbf{Rank 2 (rejects):} the remaining places are filled from this generation's new candidates (offspring,
  injected solutions, immigrants) by non-dominated sorting in (key, $V$) with crowding distance~\citep{deb2002nsga2};
  the previous generation's Rank~2 does not survive. This makes the protection of Rank~1 matter: with $y = 0$ the
  best feasible solutions are not kept by rule. A shortfall is filled with immigrants.
\end{enumerate}
Before any feasible solution is known, all candidates are ranked by $V$ (a feasibility phase). Parents are chosen by
tournaments that prefer Rank~1 to Rank~0 to Rank~2, then the lower key, then the lower violation.

\subsection{Variation and repair}
Crossover (probability 0.9) uses by default a mix of \emph{agreement} (weight 0.6), uniform (0.2) and constraint-block
(0.2) crossover. In agreement crossover the child keeps the genes on which both parents agree, and repair works only on
rows that contain a disagreeing gene --- related to the combination heuristics in MIP
solvers, which fix the variables on which pool solutions agree and solve the rest as a
sub-MIP~\citep{berthold2006primal}. Mutation changes an absolute
number of integer genes per child (Poisson, mean 4) by flips, steps, moves toward the LP solution, moves inside violated
rows or moves that improve the objective. Children are then repaired by a feasibility-jump local
search~\citep{luteberget2023feasibility}: a weighted violation over the rows, one variable at a time moved to a
\emph{jump value} that minimises it, and row weights that grow at local minima. Once an incumbent exists, an extra
row $c^\top x \le k^* - \delta$ joins the weighted rows, so that repair keeps an objective better than the incumbent
instead of trading it away. The GPU implementation runs one thread block per individual with incremental updates of the
row activities; a single-core CPU implementation of the same algorithm is the reference and the control.
Singleton continuous columns are decoded in closed form; other continuous variables are moved by their jump values.

\subsection{Exchange with the solver}
Solver and EA run as separate processes released together by a start barrier and share one monotonic clock. They
exchange solutions through a Redis server on the local machine or, as a fallback, through a mailbox directory in shared
memory; the choice is fixed by a pre-registered overhead measurement. Into SCIP, solutions enter only through a
heuristic plug-in that tries each posted solution; into HiGHS, through its user-solution callback. The solver posts its
incumbent, up to eight pool solutions within 1\,\% of it, its proven bound and, once, its root LP solution; the EA posts
a feasible solution whenever it beats the solver's incumbent. cuOpt never receives solutions (accepting them disables
its presolve); it is a control arm only.

\subsection{The frontier}
Beside the ranks the EA keeps a \emph{frontier}: an archive of feasible solutions within 5\,\% of the incumbent that
are non-dominated in (objective, distance from the incumbent), thinned to 32 by crowding distance. It is an output (a
set of diverse alternative solutions, reported with its size, spread and hypervolume), an optional source of
immigrants, and in one ablation a source of additional solutions posted to the solver. With several objectives the same
archive becomes their Pareto front, as in the first I-TREA.
